% ---------------------------------------------------------------------------
% Chapter: Examples, recording, benchmarks and development workflow
% ---------------------------------------------------------------------------
\chapter{Examples, Benchmarks and Development}\label{chap:examples}
% Long monospaced identifiers leave little room for line breaking; allow some
% extra inter-word stretch in this chapter (the group keeps it local).
\begingroup\setlength{\emergencystretch}{3em}% chapter-local

This chapter is a practical companion to the reference chapters. It lists the
example scripts that ship with the repository and the \code{env-lib} command
line, shows how to write evaluation loops and record episodes as GIF or MP4
files, explains the benchmark tools, and describes the testing, linting and
continuous-integration workflow for contributors.

\section{Example scripts}\label{sec:examples-catalogue}

The directory \code{examples/} contains runnable, documented scripts, grouped
by topic: \code{getting\_started/} (Table~\ref{tab:examples-start}),
\code{environments/} (Table~\ref{tab:examples-envs}), \code{algorithms/}
(Table~\ref{tab:examples-marl}) and \code{toolkit/}
(Table~\ref{tab:examples-toolkit}); \code{examples/README.md} lists them with a
typical invocation each. Every script has a \code{-{}-help} option, a module
docstring that explains what it demonstrates, and a \code{main()} function. Media files are written to
\code{renders/} (relative to the working directory), which Git ignores. Run
the scripts from the repository root after installing the package with the
extras they need (\code{pip install ".[all]"} covers all of them).

\begin{table}[htbp]
  \centering
  \small
  \caption{First steps (\code{examples/getting\_started/}).}
  \label{tab:examples-start}
  \begin{tabularx}{\textwidth}{@{}>{\raggedright\arraybackslash}p{3.9cm}>{\raggedright\arraybackslash}X@{}}
    \toprule
    Script & What it shows \\
    \midrule
    \code{quickstart.py} & Creates every registered environment with \code{env\_lib.make}, takes a few random steps and prints the observation shape and return; \code{-{}-render DIR} saves one PNG frame per environment. \\
    \code{workflow\_demo.py} & The workflow of Chapter~\ref{chap:workflow} in six steps: lists the environments with continuous actions, creates a native vector environment, evaluates random actions and the baseline on all copies in parallel, flattens the spaces for single-agent libraries, adapts the environment to the PettingZoo Parallel API and records a GIF of the baseline (\code{-{}-env}, \code{-{}-num-envs}, \code{-{}-episodes}, \code{-{}-save}, \code{-{}-no-save}). \\
    \bottomrule
  \end{tabularx}
\end{table}

\begin{table}[htbp]
  \centering
  \small
  \caption{Environment examples (\code{examples/environments/}).}
  \label{tab:examples-envs}
  \begin{tabularx}{\textwidth}{@{}>{\raggedright\arraybackslash}p{3.9cm}>{\raggedright\arraybackslash}X@{}}
    \toprule
    Script & What it shows \\
    \midrule
    \code{kuramoto\_demo.py} & Drives an oscillator network to synchronisation with a phase-alignment controller and a coupling ramp; NumPy or batched PyTorch backend (\code{-{}-backend}), several topologies. \\
    \code{linemsg\_demo.py} & Compares random, duty-cycled and always-relay policies on LineMsg (\code{-{}-policy} selects the one to show or record). \\
    \code{wireless\_comm\_demo.py} & Compares random access, slotted ALOHA and a collision-free rotating schedule on WirelessComm. \\
    \code{pistonball\_demo.py} & Compares a hand-written heuristic with random actions; \code{-{}-manual} enables keyboard control in a window. \\
    \code{consensus\_demo.py} & Compares random actions with the distributed Laplacian controller for consensus or formation control on any topology. \\
    \code{power\_grid\_demo.py} & Compares no control, random injections and decentralised droop control on \envid{PowerGrid-v0} under the same disturbances; reports the return, frequency nadir, largest deviation, settling time and trips (\code{-{}-buses}, \code{-{}-topology}, \code{-{}-gain}, \code{-{}-noise}). \\
    \code{platoon\_demo.py} & Compares random actions, sensor-only ACC and cooperative ACC on \envid{Platoon-v0} and prints the analytic string-stability gain of both controllers; reports collisions, the smallest gap and the peak spacing errors of the first and last follower (\code{-{}-vehicles}, \code{-{}-topology}, \code{-{}-scenario}, \code{-{}-headway}). \\
    \code{ajlatt\_demo.py} & Compares random actions with the packaged encircling baseline (\code{env\_lib.baseline\_policy}, implemented by \code{env\_lib.baselines.ajlatt\_encircle}, computed from the observation alone) for multi-robot target tracking and reports tracking error, covariance and collisions. \\
    \bottomrule
  \end{tabularx}
\end{table}

\begin{table}[htbp]
  \centering
  \small
  \caption{Toolkit examples (\code{examples/toolkit/}).}
  \label{tab:examples-toolkit}
  \begin{tabularx}{\textwidth}{@{}>{\raggedright\arraybackslash}p{5.5cm}>{\raggedright\arraybackslash}X@{}}
    \toprule
    Script & What it shows \\
    \midrule
    \code{plotkit\_gallery.py} & The three-panel paper figure of Figure~\ref{fig:plotkit-gallery}: learning curves with confidence bands, grouped bars and a sweep heatmap, saved as PDF and PNG. \\
    \code{parakit\_demo.py} & Saving, loading and applying \pkg{argparse} parameters with validation callbacks; \code{-{}-gui} opens the Tk editor, \code{-{}-load} reads a file or directory, \code{-{}-set KEY=VALUE} overrides one value. \\
    \code{simple\_transformer\_example.py} & A \code{TransformerPolicyNetwork} trained with REINFORCE on a small memory task (requires PyTorch). \\
    \code{transformer\_rl\_example.py} & PPO with transformer actor and critic networks on batched environments (requires PyTorch; \code{-{}-device auto} uses a GPU when available). \\
    \bottomrule
  \end{tabularx}
\end{table}

\begin{table}[htbp]
  \centering
  \small
  \caption{Learning examples (\code{examples/algorithms/}; require PyTorch).}
  \label{tab:examples-marl}
  \begin{tabularx}{\textwidth}{@{}>{\raggedright\arraybackslash}p{4.3cm}>{\raggedright\arraybackslash}X@{}}
    \toprule
    Script & What it shows \\
    \midrule
    \code{marl\_training\_demo.py} & \pkg{marl\_algorithms} on the environments (Chapter~\ref{chap:marl}): MAPPO on \envid{PowerGrid-v0}, MADDPG on \envid{Consensus-v0}, QMIX on \envid{LineMsg-v0} and VDN on \envid{WirelessComm-v1}, each trained with its tuned preset on batched copies and then compared with random actions and the baseline controller on the same seeded episodes; prints a results table and saves the learning curves (\code{-{}-algo}, \code{-{}-env}, \code{-{}-quick}, \code{-{}-episodes}, \code{-{}-plot}, \code{-{}-gif}, \code{-{}-seed}, \code{-{}-threads}). \\
    \code{baseline\_comparison.py} & Your method against the baselines
      (Section~\ref{sec:marl-baselines}): a distributed droop controller that
      also uses the neighbours' mean frequency deviation, compared on
      \envid{PowerGrid-v0} with MAPPO and IPPO trained over two seeds, the
      classical droop controller and random actions, on the same episodes with
      \code{marl\_algorithms.compare}; prints the table and saves it as CSV
      and a bar chart (\code{-{}-algos}, \code{-{}-seeds}, \code{-{}-episodes},
      \code{-{}-quick}, \code{-{}-csv}, \code{-{}-plot}, \code{-{}-threads}). \\
    \bottomrule
  \end{tabularx}
\end{table}

The environment examples share a set of options: \code{-{}-seed}, \code{-{}-steps}
(the episode length), \code{-{}-episodes} (evaluation episodes per policy),
\code{-{}-save PATH} (record an episode; the suffix \code{.gif} or \code{.mp4}
selects the format; \code{power\_grid\_demo.py} and \code{platoon\_demo.py} call
it \code{-{}-render PATH}, and the former also accepts \code{-{}-save}),
\code{-{}-render-mode} (\code{human} opens a window) and
\code{-{}-theme dark|light}. \code{workflow\_demo.py} takes an environment id
(\code{-{}-env}) instead, and \code{quickstart.py} only \code{-{}-steps} and
\code{-{}-render DIR}. The transformer examples,
\code{marl\_training\_demo.py} and \code{baseline\_comparison.py} accept
\code{-{}-quick} for a run of a few seconds, used as a smoke test in
continuous integration.

\begin{shellcode}
python examples/getting_started/quickstart.py --render renders/quickstart
python examples/environments/kuramoto_demo.py --topology ring --save renders/kuramoto.gif
python examples/environments/linemsg_demo.py --policy relay --save renders/linemsg.gif
python examples/environments/wireless_comm_demo.py --policy schedule --save renders/wireless.gif
python examples/environments/pistonball_demo.py --episodes 20 --save renders/pistonball.gif
python examples/environments/consensus_demo.py --task formation --save renders/formation.gif
python examples/environments/power_grid_demo.py --topology ring --render renders/grid.gif
python examples/environments/platoon_demo.py --scenario stop_and_go --render renders/platoon.gif
python examples/environments/ajlatt_demo.py --map obstacles05 --save renders/ajlatt.gif
python examples/getting_started/workflow_demo.py --env Consensus-v0 --num-envs 128 --episodes 128
python examples/toolkit/plotkit_gallery.py --formats pdf png
python examples/toolkit/simple_transformer_example.py --quick
python examples/algorithms/marl_training_demo.py --algo qmix --env LineMsg-v0
\end{shellcode}

The full suite of \code{marl\_training\_demo.py} (without \code{-{}-algo})
trains four presets and takes a few minutes on one CPU core; its results are in
Table~\ref{tab:marl-demo-results} and Figure~\ref{fig:marl-training}. With
\code{-{}-quick} every run takes only a few updates, so the trained returns
only show that the pipeline works (times and returns depend on the machine):

\begin{textcode}
Training MAPPO on PowerGrid-v0 ...
  2,048 env steps in 4.4 s; trained return -63.31
Training MADDPG on Consensus-v0 ...
  1,504 env steps in 2.3 s; trained return -1.278e+04
Training QMIX on LineMsg-v0 ...
  1,504 env steps in 1.9 s; trained return 82.17
Training VDN on WirelessComm-v1 ...
  1,504 env steps in 1.7 s; trained return 151.6
\end{textcode}

followed by a table with the columns of Table~\ref{tab:marl-demo-results}
(random $-629.3$, trained $-63.31$ and baseline $-0.7858$ for MAPPO in this
quick run) and the path of the learning-curve figure.

Table~\ref{tab:examples-results} lists the results that the comparison
scripts print with their default arguments. They are useful as sanity checks
after modifying an environment and as simple baselines; the environment
chapters explain the policies.

\begin{table}[htbp]
  \centering
  \small
  \caption{Output of the comparison scripts with default arguments (mean
    episode return, with the standard deviation over episodes where the script
    prints one).}
  \label{tab:examples-results}
  \begin{tabularx}{\textwidth}{@{}>{\raggedright\arraybackslash}p{4.3cm}>{\raggedright\arraybackslash}X@{}}
    \toprule
    Script (setting) & Result \\
    \midrule
    \code{linemsg\_demo.py} (10 agents, 20 episodes) &
      random $44.1 \pm 4.1$, duty-cycled $85.2 \pm 3.0$, always relay $95.0$ \\
    \code{wireless\_comm\_demo.py} ($6 \times 6$, 20 episodes) &
      random $385 \pm 12$, ALOHA $398 \pm 16$, rotating schedule $911 \pm 17$ with no collisions \\
    \code{pistonball\_demo.py} (20 pistons, 5 episodes) &
      heuristic $889 \pm 22$ (all episodes reach the goal), random $-342 \pm 40$ (none) \\
    \code{consensus\_demo.py} (ring, 8 agents) &
      Laplacian controller solves the task after 99 steps (return $-1968$); random actions do not (return $-19906$) \\
    \code{power\_grid\_demo.py} (16 buses, 3 episodes) &
      droop $-0.98 \pm 0.10$ (nadir $-0.019$\,Hz, settled after 2.95\,s); no control $-103.86 \pm 13.64$ (nadir $-0.114$\,Hz, never settled); random $-1020.10 \pm 726.37$ with one trip \\
    \code{platoon\_demo.py} (8 followers, 5 episodes) &
      CACC $-62.1$ without collisions over all 600 steps; ACC $-824.6$ and random $-4824.8$, both colliding in every episode; string-stability gain 1.000 (CACC) and 1.377 (ACC) \\
    \code{ajlatt\_demo.py} (\code{obstacles04}, 2 episodes) &
      baseline $-7333$ and $-7113$ with no collisions and target covariance trace about 2; random actions $-34171$ and $-20871$ with 63 and 53 collisions \\
    \code{workflow\_demo.py} (\envid{Formation-v0}, 32 copies, 64 episodes) &
      baseline $-1608.2 \pm 635.7$, all episodes successful after 92.7 steps on average; random $-18076.8 \pm 5566.8$, none successful \\
    \bottomrule
  \end{tabularx}
\end{table}

\section{The \texorpdfstring{\code{env-lib}}{env-lib} command line}\label{sec:examples-cli}

The console script \code{env-lib} (equivalently \code{python -m env\_lib}) is
the zero-code way to use every environment: it lists and describes them, runs
and records an episode of the baseline controller or of random actions,
evaluates both with confidence intervals, and benchmarks the single, native
vector and \code{SyncVectorEnv} implementations. Constructor arguments are
passed after \code{-{}-kwarg} as \code{key=value} pairs.
Section~\ref{sec:workflow-cli} documents all sub-commands and options.

\begin{shellcode}
env-lib list --continuous                  # continuous-control ids
env-lib describe Platoon-v0                # spaces, layout, parameters, baseline
env-lib run PowerGrid-v0                   # one episode of the baseline
env-lib run Platoon-v0 --kwarg scenario=stop_and_go --steps 300 \
    --theme light --gif renders/platoon.gif  # record a GIF
env-lib evaluate PowerGrid-v0 --episodes 32 --num-envs 32
env-lib bench Platoon-v0 --num-envs 256 --steps 200
\end{shellcode}

The recording command prints a summary of the episode (agent rows shortened
here):

\begin{textcode}
Platoon-v0 (scenario='stop_and_go')
  policy        baseline: cooperative adaptive cruise control (env_lib.platoon_env.cacc_policy)
  seed          0
  episode       300 steps, stopped (--steps) (29.17 s)
  team return   -26.8019
  agent return  mean -3.3502, min -3.9278 (agent_0), max -2.8442 (agent_7)
    agent     return
    agent_0  -3.9278
    ...
    agent_7  -2.8442
  final info    leader_speed=18.9  leader_acceleration=0.0009988  min_gap=11.41  collision=False  breakup=False  error_amplification=0.338  step=300  time=30
  gif           renders/platoon.gif (301 frames)
\end{textcode}

and the evaluation compares random actions with the baseline on 32 copies
simulated in parallel:

\begin{textcode}
PowerGrid-v0: 32 episodes per policy, seed 0, 32 parallel copies
policy    mean return       std  95% CI +-  length  terminated  success     s
random      -537.1770  404.5420   145.8530   198.1          6%        -  0.22
baseline      -0.8164    0.6262     0.2258   200.0          0%        -  0.15
\end{textcode}

The last column is the wall-clock time of each evaluation in seconds.

\section{Writing an evaluation loop}\label{sec:examples-loop}

All environments share the Gymnasium interface, so one evaluation function
serves all of them. The packaged \code{env\_lib.evaluate}
(Section~\ref{sec:workflow-evaluation}) does this with confidence intervals
and runs vector environments in parallel; writing the loop yourself is useful
when an experiment needs other statistics. The loop below evaluates a policy
over several seeded episodes and reports the team return and the per-agent
returns from \code{info["agent\_rewards"]}. \code{np.sum} and \code{np.any}
make it work also for AJLATT, whose reward and \code{terminated} are
per-robot arrays.

\begin{pycode}
import numpy as np
import env_lib

def evaluate(env, policy, episodes=5, seed=0):
    """Mean team return and mean per-agent returns over seeded episodes."""
    team, per_agent = [], []
    for episode in range(episodes):
        obs, info = env.reset(seed=seed + episode)
        total, agents = 0.0, 0.0
        while True:
            obs, reward, terminated, truncated, info = env.step(policy(obs))
            total += float(np.sum(reward))
            agents = agents + info["agent_rewards"]
            if np.any(terminated) or truncated:
                break
        team.append(total)
        per_agent.append(agents)
    return np.mean(team), np.mean(per_agent, axis=0)

env = env_lib.make("WirelessComm-v1")
env.action_space.seed(0)                         # reproducible random actions
mean_return, agent_returns = evaluate(env, lambda obs: env.action_space.sample())
print(f"{mean_return:.1f}", agent_returns.shape)   # 16 agents on the 4x4 grid
env.close()
\end{pycode}

\code{env\_lib.baseline\_policy(env)} returns the baseline controller of any
environment as a function of the observation, which can be passed to this
loop directly. Policies that need more than the observation can close over the
environment, for example \code{lambda obs: env.unwrapped.laplacian\_policy()}
for the consensus environment. Use \code{env.unwrapped} to reach
environment-specific methods and attributes through the wrappers that
\code{gymnasium.make} adds.

\section{Recording episodes}\label{sec:examples-recording}

The module \pkg{env\_lib.utils} provides two functions to turn rendered
frames into animations:

% norun
\begin{pycode}
record_episode(env, policy=None, path=None, *, max_steps=None, seed=None,
               fps=None, render_every=1) -> list[np.ndarray]
save_animation(frames, path, fps=30.0, *, loop=0) -> Path
\end{pycode}

\code{record\_episode} requires an environment created with
\code{render\_mode="rgb\_array"}. It resets the environment with \code{seed},
renders the initial state, then steps with \code{policy(observation)} (uniform
random actions from \code{env.action\_space} when no policy is given) until the
episode terminates or is truncated, or \code{max\_steps} steps have been
taken, rendering every \code{render\_every}-th step. It returns the list of
\code{(H, W, 3)} \code{uint8} frames and, if \code{path} is given, writes them
with \code{save\_animation} at \code{fps} frames per second (default: the
environment's \code{metadata["render\_fps"]}).

\code{save\_animation} chooses the format from the file suffix: \code{.gif} is
written with Pillow, which is always available; \code{.mp4}, \code{.webm},
\code{.avi}, \code{.mov} and \code{.mkv} require the \code{video} extra
(\pkg{imageio} and \pkg{imageio-ffmpeg}). For video formats the frames are
cropped to even width and height, which most codecs require. Parent
directories are created. \code{loop} is the GIF loop count (0 repeats
forever).

\begin{pycode}
import env_lib
from env_lib.utils import record_episode, save_animation, set_theme

set_theme("light")          # renderers created from now on use the light theme
env = env_lib.ConsensusEnv(task="formation", formation_shape="wedge",
                           render_mode="rgb_array")
frames = record_episode(env, policy=lambda obs: env.laplacian_policy(),
                        path="renders/formation.gif", seed=0, render_every=2)
save_animation(frames, "renders/formation.mp4", fps=15)   # needs the video extra
env.close()
set_theme("dark")                   # back to the default theme
\end{pycode}

\code{set\_theme("dark"|"light")} selects the colour theme of all renderers
created afterwards: the dark theme (default) suits screens and slides, the
light theme suits papers and printed material. Custom themes can be added with
\code{register\_theme}.

For full control, for example to record only part of an episode or to
annotate frames, collect the frames yourself with \code{env.render()}:

\begin{pycode}
import env_lib
from env_lib.utils import save_animation

env = env_lib.make("KuramotoOscillator-v0", render_mode="rgb_array")
obs, info = env.reset(seed=1)
frames = []
for step in range(300):
    obs, reward, terminated, truncated, info = env.step(env.action_space.sample())
    if step % 5 == 0:                       # every fifth step
        frames.append(env.render())
    if terminated or truncated:
        break
save_animation(frames, "renders/kuramoto.gif", fps=20)
env.close()
\end{pycode}

GIF files grow quickly with the number of frames and the frame size;
rendering every second or fifth step and using MP4 for long episodes keeps the
files small. The first \code{render()} call of an environment builds its
figure and takes noticeably longer than later calls.

\section{Benchmarks}\label{sec:examples-benchmarks}

\paragraph{Single environments.} The script \code{benchmarks/benchmark\_envs.py} measures the mean wall-clock
time of \code{env.step} (with pre-sampled random actions and no rendering)
and the additional time of \code{env.render()} in \code{"rgb\_array"} mode,
for representative configurations of every environment. Environments whose
optional dependencies are missing are reported as skipped. The script sets
\code{MPLBACKEND=Agg} and \code{SDL\_VIDEODRIVER=dummy} itself, so it runs
headless.

\begin{shellcode}
python benchmarks/benchmark_envs.py                        # all environments
python benchmarks/benchmark_envs.py --steps 500 --only kuramoto ajlatt
python benchmarks/benchmark_envs.py --markdown > renders/benchmarks.md
\end{shellcode}

\begin{table}[htbp]
  \centering
  \small
  \caption{Options of \code{benchmarks/benchmark\_envs.py}.}
  \label{tab:examples-bench-options}
  \begin{tabularx}{\textwidth}{@{}l>{\raggedright\arraybackslash}p{1.6cm}>{\raggedright\arraybackslash}X@{}}
    \toprule
    Option & Default & Meaning \\
    \midrule
    \code{-{}-steps} & \code{300} & Timed steps per case. \\
    \code{-{}-render-steps} & \code{40} & Timed rendered steps per case; \code{0} skips the rendering measurement. \\
    \code{-{}-seed} & \code{0} & Seed of the environment and of the action space. \\
    \code{-{}-only} & all & Subset of the groups \code{kuramoto}, \code{linemsg}, \code{wireless}, \code{pistonball}, \code{consensus}, \code{ajlatt}, \code{power\_grid}, \code{platoon}. \\
    \code{-{}-markdown} & off & Print a Markdown table instead of plain text. \\
    \code{-{}-torch-threads} & PyTorch default & Number of PyTorch threads for the PyTorch Kuramoto case. \\
    \bottomrule
  \end{tabularx}
\end{table}

Table~\ref{tab:examples-bench} shows typical results. The ``before 1.0''
column was measured with the implementations that preceded version 1.0; the
version 1.1 columns are the output of
\code{OMP\_NUM\_THREADS=1 python benchmarks/benchmark\_envs.py -{}-torch-threads 1}
on one core of an Intel Xeon at 2.8~GHz (Python~3.11, NumPy~2.4). Timings on a
shared machine vary by 10 to 20 per cent between runs; the ratios are more
informative than the absolute numbers.

\begin{table}[htbp]
  \centering
  \small
  \caption{Time per step and per rendered \code{"rgb\_array"} frame in
    milliseconds. A dash marks a configuration that did not exist before
    version 1.0.}
  \label{tab:examples-bench}
  \begin{tabular}{@{}lrrr@{}}
    \toprule
    Configuration & Step, before 1.0 & Step, 1.1 & Frame, 1.1 \\
    \midrule
    Kuramoto NumPy, $N = 10$              & 0.17  & 0.075 & 28 \\
    Kuramoto NumPy, $N = 50$              & 2.36  & 0.093 & 44 \\
    Kuramoto NumPy RK4, $N = 50$          & 7.1   & 0.12  & 53 \\
    Kuramoto PyTorch, $N = 50$, 8 systems & 10.3  & 0.45  & 42 \\
    LineMsg, 10 agents                    & 0.019 & 0.017 & 16 \\
    WirelessComm, $6 \times 6$            & 0.154 & 0.057 & 14 \\
    WirelessComm, $12 \times 12$          & 0.56  & 0.064 & 18 \\
    Pistonball, 20 pistons                & 1.11  & 0.078 & 3 \\
    Consensus, 8 agents                   & --    & 0.082 & 19 \\
    Formation, 32 agents, proximity graph & --    & 0.28  & 29 \\
    PowerGrid, 16 buses                   & --    & 0.25  & 36 \\
    Platoon, 8 followers                  & --    & 0.13  & 38 \\
    AJLATT, \code{obstacles04}, 4 robots  & 90    & 4.7   & 33 \\
    AJLATT, \code{obstacles05}, 4 robots  & 67    & 4.7   & 34 \\
    \bottomrule
  \end{tabular}
\end{table}

Before version 1.0 the Kuramoto \code{"rgb\_array"} renderer failed with
matplotlib 3.10 and later, Pistonball drew its screen on every step even
without rendering, and an AJLATT frame took about 110\,ms. Rendering now
costs more than a simulation step in every environment, which is why it is
performed only on request (Section~\ref{sec:trouble-performance}). The
chapters of the environments added in version 1.1 report their step and frame
times (Sections~\ref{sec:power-grid-vector}, \ref{sec:power-grid-render},
\ref{sec:platoon-vector} and~\ref{sec:platoon-render}).

\paragraph{Vector environments.} \code{benchmarks/benchmark\_vector.py} measures
the throughput of \code{env\_lib.make\_vec(id, n)}, once with the native
implementation and once with Gymnasium's \code{SyncVectorEnv}, for
$n \in \{1, 16, 64, 256, 1024\}$
(Sync is skipped above \code{-{}-sync-max}, default 256) and prints environment
steps, agent-steps per second and the speed-up; \code{-{}-ids} selects
environments, \code{-{}-num-envs} the batch sizes, \code{-{}-markdown} prints a
Markdown table and \code{-{}-quick} runs a short smoke test. Every case runs at
least 300 batched steps (\code{-{}-min-steps}), so automatic resets are
included. \code{env-lib bench ID} measures the throughput
of an environment in environment steps per second: the single environment, its
native vector implementation (if any) and \code{SyncVectorEnv} over single
environments (and \code{AsyncVectorEnv} with \code{-{}-async}), all stepped
with a fixed random action. \code{-{}-num-envs} sets the number of copies,
\code{-{}-steps} the number of timed steps and \code{-{}-kwarg} the
constructor arguments. Table~\ref{tab:workflow-throughput} collects reference
values, and Section~\ref{sec:trouble-performance} discusses batch sizes and
threads.

\begin{shellcode}
env-lib bench Consensus-v0 --num-envs 256 --steps 200
env-lib bench PowerGrid-v0 --num-envs 1024 --kwarg n_buses=64
env-lib bench LineMsg-v0 --num-envs 64 --async     # no native implementation
\end{shellcode}

\paragraph{Algorithms.} \code{benchmarks/benchmark\_marl.py} trains every
preset of \pkg{marl\_algorithms} (seed 0, one thread) and evaluates random
actions, the trained policy and the baseline controller on the same 64 seeded
episodes; \code{-{}-only} selects presets, \code{-{}-shard K/N} splits the
presets between parallel processes, \code{-{}-out} names the JSON-lines result
file and \code{-{}-report FILE} (with \code{-{}-markdown}) prints the table.
Section~\ref{sec:marl-results} describes the protocol and reports the results.

\section{Testing, linting and continuous integration}\label{sec:examples-contributing}

\subsection{Running the tests}

The test suite lives in \code{tests/} (\code{tests/envs/} for the
environments, \code{tests/algorithms/} for \pkg{marl\_algorithms},
\code{tests/toolkit/} for the toolkit) and uses \pkg{pytest}.
After an editable installation with the \code{dev} extra
(Section~\ref{sec:install-dev}):

\begin{shellcode}
pytest -q                                   # whole suite (about two minutes)
pytest tests/envs/test_consensus.py -q      # one file
pytest -k "seed or determinism" -q          # tests whose name matches
pytest --cov=env_lib --cov=marl_algorithms --cov=toolkit \
    --cov-report=term-missing
\end{shellcode}

The suite contains more than 1300 tests, 176 of them for
\pkg{marl\_algorithms}. \code{tests/conftest.py} sets
\code{MPLBACKEND=Agg} and \code{SDL\_VIDEODRIVER=dummy}, so rendering tests
run without a display. Tests of optional features are skipped when the
dependency is missing (\code{pytest.importorskip}). The marker \code{slow} is
reserved for long-running tests and can be deselected with
\code{-m "not slow"}.

The tests check, for every environment, the shapes, dtypes and bounds of the
spaces, that observations lie in \code{observation\_space}, determinism under
seeding, reward and termination rules against simple reference
implementations written in the tests, backward-compatible arguments, the
registered ids, rendered frames and, where the multi-agent interface permits
it, Gymnasium's \code{check\_env}. For the native vector environments they
check that copy~0 reproduces a single environment started from the same state,
the three autoreset modes and the info masks; for the workflow layer they
check the catalogue, the command line, the baseline controllers, the wrappers
(including PettingZoo's \code{parallel\_api\_test} when PettingZoo is
installed) and the evaluation statistics. For \pkg{marl\_algorithms} they check the per-agent
view of every kind of action space, the runner's final observations and
per-agent rewards, generalised advantage estimation against the recursion
written out, the replay buffer, the running statistics and the networks
(including the monotonicity of the QMIX mixer), and for every algorithm the
action shapes and bounds, \code{policy()} with \code{env\_lib.evaluate},
the rejection of unsupported action kinds and invalid configurations, short
training runs with shared and separate parameters, the losses and targets
against their definitions, seed determinism, \code{save}/\code{load} round
trips, small learning checks, the presets and the \code{marl-train} command
line. For the toolkit they cover input handling and numerical
results of \pkg{plotkit}, network shapes and the version 1.0 fixes of
\pkg{neural\_toolkit}, and the conversion rules of \pkg{parakit}.

\subsection{Linting and formatting}

The code base is formatted and linted with \pkg{ruff}; the configuration
(line length 100, target Python 3.9, rule sets \code{E}, \code{F}, \code{W},
\code{I}, \code{B}, \code{UP}, \code{NPY}) is in \code{pyproject.toml}.

\begin{shellcode}
ruff check src tests examples benchmarks          # lint
ruff check --fix src tests examples benchmarks    # apply automatic fixes
ruff format src tests examples benchmarks         # format
ruff format --check src tests examples benchmarks # verify formatting only
\end{shellcode}

\subsection{Continuous integration}

The workflow \code{.github/workflows/ci.yml} runs on pushes to the
\code{main} and \code{development} branches and on pull requests. It has two
jobs:

\begin{description}[leftmargin=0.6cm,style=nextline]
  \item[\code{lint}] On Python 3.11: \code{ruff check} and
    \code{ruff format -{}-check} on \code{src}, \code{tests}, \code{examples} and
    \code{benchmarks}.
  \item[\code{test}] On Python 3.9, 3.10, 3.11 and 3.12 with
    \code{MPLBACKEND=Agg} and \code{SDL\_VIDEODRIVER=dummy}: installs the CPU
    build of PyTorch and the package with the \code{pistonball}, \code{video}
    and \code{dev} extras, runs \code{pytest -q}, and smoke-tests the examples
    and the command line: \code{quickstart.py -{}-steps 5},
    \code{simple\_transformer\_example.py -{}-quick -{}-no-plot},
    \code{workflow\_demo.py -{}-no-save},
    \code{power\_grid\_demo.py -{}-episodes 2 -{}-steps 50},
    \code{platoon\_demo.py -{}-episodes 2 -{}-steps 100}, \code{env-lib list},
    \code{env-lib evaluate Formation-v0 -{}-episodes 8 -{}-num-envs 8},
    \code{benchmarks/benchmark\_envs.py -{}-steps 20 -{}-render-steps 2},
    \code{benchmarks/benchmark\_vector.py -{}-quick}, \code{marl-train list},
    \code{marl\_training\_demo.py -{}-quick -{}-plot ''},
    \code{baseline\_comparison.py -{}-quick} (one seed, 16 episodes) and
    \code{marl-train compare LineMsg-v0} with IQL and VDN on a small budget.
\end{description}

Running the same commands locally before opening a pull request avoids
failed CI runs.

\subsection{Conventions for contributions}

Contributions follow the conventions that the existing code relies on:

\begin{itemize}
  \item Python 3.9 compatible code with \code{from \_\_future\_\_ import
    annotations}, type hints on public functions and NumPy-style docstrings
    (Parameters, Returns, Raises, Examples).
  \item No global side effects in library code: no \code{print}, no
    \code{logging.basicConfig}, no calls to the global NumPy or PyTorch random
    generators, no changes to \code{matplotlib.rcParams} and no
    \code{matplotlib.use}. Randomness comes from \code{self.np\_random} or an
    explicitly owned generator.
  \item Arguments are validated early with \code{ValueError} or
    \code{TypeError} and actionable messages; user-facing problems are
    reported with \code{warnings.warn(..., stacklevel=2)}; diagnostics go to
    \code{logging.getLogger(\_\_name\_\_)}.
  \item Environments follow the contract of Chapter~\ref{chap:envlib}:
    Gymnasium API, \code{float32} observations inside
    \code{observation\_space}, \code{info["agent\_rewards"]}, their own episode
    limit, and no drawing work unless a render mode is set. Renderers subclass
    \code{env\_lib.utils.MatplotlibRenderer}, build their artists once and use
    the colours of the active theme.
  \item Deprecated arguments keep working and emit a
    \code{DeprecationWarning} that names the replacement; every behavioural
    change is recorded in \code{CHANGELOG.md}.
  \item Plain ASCII text in code, output and documentation, without emojis or
    decorative symbols.
\end{itemize}

A new environment needs the following additions:
\begin{enumerate}
  \item an entry in \code{ENV\_SPECS} in \code{src/env\_lib/registration.py},
    with a string entry point so that registration does not import optional
    dependencies, and the metadata of Table~\ref{tab:envlib-spec}
    (\code{family}, observation and action types, \code{requires},
    \code{limit\_kwarg} and, for a native batched implementation,
    \code{vector\_entry\_point}); the catalogue and the command line pick the
    entry up automatically;
  \item a lazy export in \code{src/env\_lib/\_\_init\_\_.py};
  \item a baseline controller that is a pure function of the observation,
    registered in \code{src/env\_lib/baselines.py} so that
    \code{baseline\_policy} finds it;
  \item for a native vector environment, a subclass of
    \code{env\_lib.utils.BatchedVectorEnv} that accepts the constructor
    arguments of the single environment plus \code{num\_envs} and
    \code{autoreset\_mode} (Appendix~\ref{app:api});
  \item tests in \code{tests/envs/} and an example script in
    \code{examples/environments/};
  \item a chapter in this manual and an entry in \code{CHANGELOG.md}.
\end{enumerate}

Section~\ref{sec:marl-extending} lists the corresponding steps for a new
learning algorithm in \pkg{marl\_algorithms}.

\par\endgroup
